\section{An exactly solvable two-level Hamiltonian}
Fix $g>0$ and $0\leq\gamma<g$. Let $\sigma_x,\sigma_y,\sigma_z$ denote Pauli
matrices and consider
\begin{equation}\label{eq:Hmodel}
 H_{g,\gamma}=g\sigma_x+i\gamma\sigma_z
 =\begin{pmatrix}i\gamma&g\\g&-i\gamma\end{pmatrix},
 \qquad E=\sqrt{g^2-\gamma^2}.
\end{equation}
Its eigenvalues are $\pm E$ and it satisfies $H_{g,\gamma}^2=E^2\I$.
At $\gamma=g$ both eigenvalues coalesce at zero and
$H_{g,g}\ne 0$ is defective. All positive-metric statements below therefore
apply only to the open interval $\gamma<g$.

\begin{proposition}[Complete positive metric family]\label{prop:metricfamily}
Up to a positive overall scalar, every positive metric compatible with
$H_{g,\gamma}$ is
\begin{equation}\label{eq:metricfamily}
 \eta_{r,t}=\I+r\sigma_y+t\sigma_x,
 \qquad r=\frac{\gamma}{g},\qquad r^2+t^2<1,
 \quad t\in\Rn.
\end{equation}
Its eigenvalues are $1\pm\sqrt{r^2+t^2}$ and
\begin{equation}\label{eq:metriccondition}
 \kg(\eta_{r,t})
 =\frac{1+\sqrt{r^2+t^2}}{1-\sqrt{r^2+t^2}}.
\end{equation}
The unique minimizing value of $t$ is $t=0$ (for a fixed overall normalization),
and
\begin{equation}\label{eq:optimalmetric}
 \eta_*=\I+\frac{\gamma}{g}\sigma_y,
 \qquad \min_{\eta>0:\,H^\dagger\eta=\eta H}\kg(\eta)
 =\frac{g+\gamma}{g-\gamma}.
\end{equation}
\end{proposition}
\begin{proof}
Every $2\times2$ Hermitian matrix has a unique expansion
$\eta=a\I+b\sigma_x+c\sigma_y+d\sigma_z$ with real coefficients.
Substitution into $H^\dagger\eta=\eta H$ gives
$d=0$ and $c=(\gamma/g)a$. Positivity requires
$a>\sqrt{b^2+c^2}$. Rescaling by $a>0$ and setting $t=b/a$
yields~\eqref{eq:metricfamily}. Diagonalizing its Pauli-vector part
gives~\eqref{eq:metriccondition}, which is increasing in $\abs{t}$.
\end{proof}

\begin{proposition}[Exact inversion, singular values and sharp certificate]
\label{prop:modelsharp}
For the minimizing metric~\eqref{eq:optimalmetric},
\begin{align}
 h_*:=\eta_*^{1/2}H_{g,\gamma}\eta_*^{-1/2}
 &=E\sigma_x, \label{eq:hexact}\\
 H_{g,\gamma}^{-1}&=\frac{H_{g,\gamma}}{E^2},\label{eq:hinverse}\\
 \norm{H_{g,\gamma}^{-1}}_2&=\frac{1}{g-\gamma},
 &\norm{h_*^{-1}}_2&=\frac{1}{E},\label{eq:inverseexact}\\
 \kg(H_{g,\gamma})&=\kg(\eta_*)=\frac{g+\gamma}{g-\gamma},
 &\kg(h_*)&=1.\label{eq:conditionexact}
\end{align}
Moreover,
\begin{equation}\label{eq:sharp}
 r_{\mathrm{sing}}(H_{g,\gamma})=g-\gamma
 =\frac{E}{\sqrt{\kg(\eta_*)}}.
\end{equation}
Thus the sufficient perturbation radius~\eqref{eq:radiusbound} is sharp
for this example and this metric.
\end{proposition}
\begin{proof}
Using the Pauli algebra and the anticommutation of $\sigma_x$ and $\sigma_y$,
$\eta_* H=(E^2/g)\sigma_x$ and
$\eta_*^{-1/2}\sigma_x\eta_*^{-1/2}=(g/E)\sigma_x$.
These identities give~\eqref{eq:hexact}. Since $H^2=E^2\I$,
\eqref{eq:hinverse} follows. A direct multiplication gives
\begin{equation}
 H^\dagger H=(g^2+\gamma^2)\I+2g\gamma\sigma_y,
\end{equation}
whose eigenvalues are $(g+\gamma)^2$ and $(g-\gamma)^2$.
Hence the singular values are $g\pm\gamma$, proving
\eqref{eq:inverseexact},~\eqref{eq:conditionexact}, and
$r_{\mathrm{sing}}(H)=g-\gamma$. Finally,
$E^2=(g-\gamma)(g+\gamma)$ gives~\eqref{eq:sharp}.
\end{proof}

\begin{remark}[A crucial distinction near the exceptional point]
As $\gamma\uparrow g$, the Hermitian representative satisfies
$\kg(h_*)=1$ for every $\gamma<g$, while
$\norm{h_*^{-1}}_2=E^{-1}\to\infty$.
Its \emph{relative} condition number being one does not prevent
\emph{absolute} inverse amplification. In the reference representation,
$\norm{H^{-1}}_2=(g-\gamma)^{-1}$ diverges even faster.
No positive definite metric persists at the defective endpoint.
\end{remark}

\subsection{The role of eigenvector overlap}
The preceding metric computation has a useful two-dimensional geometric
interpretation. We record it as an elementary auxiliary result, without
claiming that general metric optimization is new~\cite{Krejcirik2018}.

\begin{proposition}[Two-eigenvector overlap and condition-number minimum]
\label{prop:overlap}
Let a $2\times2$ matrix have distinct real eigenvalues and Euclidean-unit
right eigenvectors $v_1,v_2$. Put $c=\abs{v_1^\dagger v_2}<1$.
Among all positive metrics making the eigenvectors orthogonal, the minimum is
\begin{equation}\label{eq:overlap}
 \min_\eta\kg(\eta)=\frac{1+c}{1-c}.
\end{equation}
It is attained, up to overall positive scale, by
$\eta_*=(VV^\dagger)^{-1}$, where $V=[v_1\ v_2]$.
\end{proposition}
\begin{proof}
Adjust the phase of $v_2$ to make $v_1^\dagger v_2=c\geq0$.
An orthonormal basis $e_+,e_-$ can be chosen so that
$v_1=a e_++b e_-$ and $v_2=a e_+-b e_-$ with
$a^2=(1+c)/2$ and $b^2=(1-c)/2$.
In this basis, $v_1^\dagger\eta v_2=0$ implies
$a^2\eta_{++}=b^2\eta_{--}$ and
$\operatorname{Im}\eta_{+-}=0$.
Since $\lambda_{\min}(\eta)\leq\eta_{++}$ and
$\lambda_{\max}(\eta)\geq\eta_{--}$, it follows that
$\kg(\eta)\geq\eta_{--}/\eta_{++}=a^2/b^2$.
The choice $(VV^\dagger)^{-1}$ is diagonal in $e_\pm$ with
entries $(1+c)^{-1},(1-c)^{-1}$ and attains the bound.
Equality forces the off-diagonal entry to vanish, giving uniqueness up to scale.
\end{proof}

For~\eqref{eq:Hmodel}, normalized right eigenvectors have overlap
$c=\gamma/g$. Thus~\eqref{eq:overlap} reproduces
\eqref{eq:optimalmetric} and makes explicit the link between
coalescing eigendirections and poor metric conditioning. The overlap
identity is a diagnostic for this model, not a claim of universal
exceptional-point scaling in arbitrary dimensions.
